\documentclass[11pt]{article}
\input{../../shared/project_style.tex}
\rhead{Basics 24 Textbook Draft}

\begin{document}
\DocTitle{Basics 24: Derivation of Ideal Magnetohydrodynamics}{IV - Magnetohydrodynamics}{From multifluid plasma dynamics and Maxwell equations to the closed ideal-MHD system}

\begin{overviewbox}
\textbf{Textbook draft, version 1.0.} Ideal MHD is a controlled long-wavelength, low-frequency reduction of plasma dynamics. This chapter derives its equations, records every major approximation, and emphasizes exactly which kinetic and nonideal effects are removed. This chapter is designed for a reader with no prior plasma-physics training. It distinguishes exact identities from physical approximations and connects the central equations to later simulation modules.
\end{overviewbox}

\ProjectTOC

\section{Learning objectives}
After completing this note, the reader should be able to:
\begin{itemize}[leftmargin=1.6em]
\item State the closed ideal-MHD equations.
\item Derive mass and momentum conservation from species equations.
\item Reduce generalized Ohm's law to the ideal condition.
\item Derive the induction equation.
\item Derive pressure and total-energy closures.
\item Construct an explicit assumption ledger for model validity.
\end{itemize}

\section{Prerequisites and reading strategy}
\begin{itemize}[leftmargin=1.6em]
\item Basics 12, 21, 22, and 23; Basics 3 and 5 for vector/PDE manipulations.
\end{itemize}
On a first reading, emphasize physical meaning and dimensional checks. On a second reading, reproduce the derivations. Attempt the computational laboratory only after the limiting cases are understood.

\section{Target ideal-MHD system}

A common conservative ideal-MHD system is
\begin{align}
\partial_t\rho+\nabla\cdot(\rho\mathbf u)&=0,\\
\partial_t(\rho\mathbf u)+\nabla\cdot\left[
\rho\mathbf u\mathbf u+
\left(p+\frac{B^2}{2\mu_0}\right)\mathsf I
-\frac{\mathbf B\mathbf B}{\mu_0}\right]&=0,\\
\partial_t\mathbf B-\nabla\times(\mathbf u\times\mathbf B)&=0,\\
\partial_tE+\nabla\cdot\left[
\left(E+p+\frac{B^2}{2\mu_0}\right)\mathbf u
-\frac{\mathbf B(\mathbf u\cdot\mathbf B)}{\mu_0}\right]&=0,\\
\nabla\cdot\mathbf B&=0,
\end{align}
with
\begin{equation}
E=\frac{\rho u^2}{2}+\frac{p}{\gamma-1}+\frac{B^2}{2\mu_0}.
\end{equation}
We now identify where each equation comes from.


\section{Mass conservation}

Summing species number equations weighted by mass gives
\begin{equation}
\partial_t\rho+\nabla\cdot(\rho\mathbf u)=0
\end{equation}
when nuclear reactions, external injection, and mass loss are absent. This step requires no quasineutral approximation; it follows from total mass accounting.


\section{Momentum and magnetic stress}

Summing species momentum gives, under quasineutrality and negligible relative-drift inertia,
\begin{equation}
\rho\frac{D\mathbf u}{Dt}=-\nabla p+\mathbf J\times\mathbf B.
\end{equation}
At low frequency, Ampere's law gives $\mathbf J=\nabla\times\mathbf B/\mu_0$. Using
\begin{equation}
(\nabla\times\mathbf B)\times\mathbf B
=-\nabla\left(\frac{B^2}{2}\right)+(\mathbf B\cdot\nabla)\mathbf B
\end{equation}
for $\nabla\cdot\mathbf B=0$,
\begin{equation}
\mathbf J\times\mathbf B
=-\nabla\left(\frac{B^2}{2\mu_0}\right)
+\frac{(\mathbf B\cdot\nabla)\mathbf B}{\mu_0}.
\end{equation}
These terms are magnetic pressure and magnetic tension. Rewriting them as the divergence of Maxwell stress yields the conservative momentum equation.


\section{Ideal Ohm's law as an ordering}

Generalized Ohm's law contains
\begin{equation}
\mathbf E+\mathbf u\times\mathbf B
=\eta\mathbf J+\frac{\mathbf J\times\mathbf B}{en_e}
-\frac{\nabla\cdot\mathsf P_e}{en_e}
-\frac{m_e}{e}\frac{D_e\mathbf u_e}{Dt}+\cdots.
\end{equation}
Ideal MHD assumes all right-hand terms are small on the scales of interest:
\begin{equation}
\boxed{\mathbf E+\mathbf u\times\mathbf B=0}.
\end{equation}
This is not a universal law. It fails at sufficiently small scales, high frequencies, low collisionality, reconnection layers, and kinetic resonances.


\section{Induction equation}

Insert ideal Ohm's law into Faraday's law:
\begin{equation}
\frac{\partial\mathbf B}{\partial t}
=-\nabla\times\mathbf E
=\nabla\times(\mathbf u\times\mathbf B).
\end{equation}
Together with $\nabla\cdot\mathbf B=0$, this transports and stretches magnetic flux with the fluid. Expanding the curl gives
\begin{equation}
\frac{D\mathbf B}{Dt}
=(\mathbf B\cdot\nabla)\mathbf u
-\mathbf B(\nabla\cdot\mathbf u),
\end{equation}
after using mass conservation and zero magnetic divergence in the appropriate form.


\section{Thermodynamic closure}

For an adiabatic ideal gas with no heat conduction or dissipation,
\begin{equation}
\frac{Dp}{Dt}+\gamma p\nabla\cdot\mathbf u=0,
\end{equation}
equivalently
\begin{equation}
\frac{D}{Dt}\left(\frac{p}{\rho^\gamma}\right)=0.
\end{equation}
The parameter $\gamma$ represents the chosen closure, commonly $5/3$ for an isotropic monatomic gas. Collisionless magnetized plasmas may not justify this scalar adiabatic law.


\section{Total energy conservation}

Dot the momentum equation with $\mathbf u$, combine the internal-energy equation with Poynting's theorem, and use ideal Ohm's law. The result is
\begin{equation}
\partial_tE+\nabla\cdot\left[
\left(E+p+\frac{B^2}{2\mu_0}\right)\mathbf u
-\frac{\mathbf B(\mathbf u\cdot\mathbf B)}{\mu_0}\right]=0.
\end{equation}
The flux contains advected kinetic/internal energy, pressure work, and electromagnetic transport. A numerical solver that updates the separate pieces inconsistently can violate this balance.


\section{Assumption ledger and domain of validity}

Ideal MHD commonly assumes:
\begin{enumerate}[leftmargin=1.8em]
\item quasineutrality and negligible displacement current;
\item scales large relative to Debye length and relevant gyro radii;
\item frequencies below electron and often ion kinetic scales;
\item sufficiently small Hall, resistive, electron-inertia, and anisotropic-pressure effects;
\item a scalar thermodynamic closure;
\item a continuum description with enough particles per resolved volume.
\end{enumerate}
The present project uses ideal MHD for bulk Alfvén-wave structure, then adds energetic-particle kinetics because resonant fast-ion drive lies outside this closure.


\section{Computational laboratory}
Create an assumption-ledger notebook that evaluates dimensionless ratios for a selected fusion-plasma case. Numerically verify the magnetic-force identity and conservative stress form on a smooth divergence-free field. Integrate a one-dimensional ideal-MHD wave benchmark and monitor total energy.

\section{Summary}
\begin{itemize}[leftmargin=1.6em]
\item Ideal MHD is derived by combining species moments, low-frequency Maxwell equations, and an ordered Ohm law.
\item Magnetic pressure and tension arise from the Lorentz force.
\item Ideal Ohm's law produces the induction equation and frozen magnetic flux.
\item A scalar adiabatic law closes the thermodynamics.
\item Every ideal-MHD result inherits a clear scale and closure domain.
\end{itemize}

\SymbolsAcronymsSection
\begin{symtable}
$\rho,\mathbf u,p$ & fluid variables & Mass density, bulk velocity, and scalar pressure. \\
$E$ & total energy density & Kinetic, internal, and magnetic energy. \\
$\gamma$ & adiabatic index & Thermodynamic closure parameter. \\
$\eta$ & resistivity & Nonideal coefficient omitted in ideal MHD. \\
$D/Dt$ & material derivative & $\partial_t+\mathbf u\cdot\nabla$. \\
$\mathsf I$ & identity tensor & Unit second-order tensor. \\
\end{symtable}

\section{Exercises}
\begin{enumerate}[leftmargin=1.8em]
\item Derive the magnetic-pressure and tension decomposition of $\mathbf J\times\mathbf B$.
\item Derive the induction equation from Faraday and ideal Ohm laws.
\item Show that $p/\rho^\gamma$ is advected for the adiabatic closure.
\item List the terms omitted from generalized Ohm's law and associate each with a physical scale.
\item Explain why ideal MHD can describe an Alfvén eigenfunction but not energetic-particle resonant drive by itself.
\end{enumerate}

\section{References}
\begin{thebibliography}{99}
\bibitem{ref1} D. J. Griffiths, \emph{Introduction to Electrodynamics}, 4th ed., Cambridge University Press (2017).
\bibitem{ref2} J. D. Jackson, \emph{Classical Electrodynamics}, 3rd ed., Wiley (1998).
\bibitem{ref3} F. F. Chen, \emph{Introduction to Plasma Physics and Controlled Fusion}, 3rd ed., Springer (2016).
\bibitem{ref4} R. J. Goldston and P. H. Rutherford, \emph{Introduction to Plasma Physics}, CRC Press (1995).
\bibitem{ref5} P. M. Bellan, \emph{Fundamentals of Plasma Physics}, Cambridge University Press (2006).
\bibitem{ref6} J. P. Freidberg, \emph{Ideal MHD}, Cambridge University Press (2014).
\bibitem{ref7} J. P. Goedbloed, R. Keppens, and S. Poedts, \emph{Advanced Magnetohydrodynamics}, Cambridge University Press (2010).
\bibitem{ref8} H. Goedbloed and S. Poedts, \emph{Principles of Magnetohydrodynamics}, Cambridge University Press (2004).
\bibitem{ref9} T. H. Stix, \emph{Waves in Plasmas}, American Institute of Physics (1992).
\bibitem{ref10} D. A. Gurnett and A. Bhattacharjee, \emph{Introduction to Plasma Physics}, 2nd ed., Cambridge University Press (2017).
\end{thebibliography}

\end{document}
